\section*{Capítulo \ref{ch:b3:radicals-carbenes} --- Radicais, carbenos e rearranjos}

\begin{solution}{exo:b3:radicals-carbenes:1}
O propano tem seis hidrogênios primários e dois secundários. Cloração: $6 \times 1 : 2 \times 3.9 = 6 : 7.8$,
isto é, 43\,\% de 1-cloropropano e 57\,\% de 2-cloropropano. Bromação: $6 : 164$, 4\,\% e
96\,\%.
\end{solution}

\begin{solution}{exo:b3:radicals-carbenes:2}
$\ce{CH2}$: tripleto; $\ce{CCl2}$: singleto (os pares isolados do cloro elevam o orbital p vazio).
O diclorocarbeno se adiciona de modo estereoespecífico: \textit{trans}-1,1-dicloro-2,3-dimetilciclopropano.
\end{solution}

\begin{solution}{exo:b3:radicals-carbenes:3}
\textit{cis}-1,2-Dietilciclopropano: o \omterm{def:b3:radicals-carbenes:carbenoid}{carbenoide} se adiciona em uma única etapa a uma das faces.
\end{solution}

\begin{solution}{exo:b3:radicals-carbenes:4}
Ciclo-hexanona: $\varepsilon$-caprolactona (oxepan-2-ona). Acetofenona: acetato de fenila,
$\ce{CH3CO2C6H5}$; a fenila migra preferencialmente à metila.
\end{solution}

\begin{solution}{exo:b3:radicals-carbenes:5}
Nove hidrogênios primários e um terciário. Cloração: $9 : 5.1$, 64\,\% de 1-cloro-2-metilpropano
e 36\,\% de 2-cloro-2-metilpropano. Bromação: $9 : 1600$, 0.6\,\% e 99.4\,\%.
\end{solution}

\begin{solution}{exo:b3:radicals-carbenes:6}
$k_c = 3.0 \times k_H[\mathrm{Bu_3SnH}] = 3.0 \times \num{2.4e6} \times 0.020 = \qty{1.4e5}{s^{-1}}$.
\end{solution}

\begin{solution}{exo:b3:radicals-carbenes:7}
Ataque sobre C5: 5-exo-trig, favorecido; sobre C6: 6-endo-trig, favorecido pelas regras, mas
mais lento aqui. 4-exo-tet: favorecido. 5-endo-trig: desfavorecido. 5-exo-dig: favorecido.
3-exo-dig: desfavorecido. 5-endo-dig: favorecido.
\end{solution}

\begin{solution}{exo:b3:radicals-carbenes:8}
A perda de água dá um cátion terciário e benzílico; no carbono vizinho, uma fenila
e uma metila poderiam migrar, e a fenila, de maior aptidão, migra. O
cátion estabilizado pelo oxigênio perde um próton: 3,3-difenilbutan-2-ona.
\end{solution}

\begin{solution}{exo:b3:radicals-carbenes:9}
Curtius: $\ce{C6H5CON3 -> C6H5NCO + N2}$. Com água, o ácido carbâmico perde $\ce{CO2}$: anilina
(forma-se um pouco de difenilureia quando a anilina encontra mais \omterm{def:b3:radicals-carbenes:curtius}{isocianato}). Com etanol:
\textit{N}-fenilcarbamato de etila, $\ce{C6H5NHCO2C2H5}$.
\end{solution}

\begin{solution}{exo:b3:radicals-carbenes:10}
Cloro: primário $421.7 - 432 = \qty{-10}{kJ/mol}$, terciário $405.4 - 432 = \qty{-27}{kJ/mol}$; bromo: $+56$ e
$+39$. As duas abstrações pelo cloro são exotérmicas, com estados de transição precoces que
sentem pouco a diferença de \qty{16}{kJ/mol}; as duas abstrações pelo bromo são endotérmicas,
com estados de transição tardios cujas energias de ativação diferem pela maior parte dela:
$\eu^{16000/(8.314 \times 298)} \approx 600$, a ordem de grandeza da razão observada.
\end{solution}

\begin{solution}{exo:b3:radicals-carbenes:11}
Oxima \textit{E} (OH anti a C2, que porta a metila): C2 migra, e o nitrogênio entra
ao lado dele: 7-metilazepan-2-ona. Oxima \textit{Z}: C6 migra: 3-metilazepan-2-ona. No
\omterm{def:b3:radicals-carbenes:beckmann}{rearranjo de Beckmann}, a geometria da oxima decide; na \omterm{def:b3:radicals-carbenes:baeyer}{oxidação de Baeyer--Villiger},
é a \omterm{def:b3:radicals-carbenes:migratory}{aptidão migratória}, e o carbono secundário migra:
7-metiloxepan-2-ona.
\end{solution}

\begin{solution}{exo:b3:radicals-carbenes:12}
Fração ciclizada $k_c/(k_c + k_H c)$: 0.91, 0.49 e 0.09. Para 95\,\%: $k_Hc = k_c(1/0.95 - 1)$, logo
$c = 0.0526 \times \num{2.3e5}/\num{2.4e6} = \qty{5.0e-3}{mol/L}$. Ela é mantida adicionando o hidreto lentamente com uma
bomba de seringa, ou usando uma quantidade catalítica de haleto de estanho com um redutor
que regenera o hidreto.
\end{solution}

\begin{solution}{pb:b3:radicals-carbenes:1}
\textbf{1.} $\ce{C6H10O + NH2OH -> C6H11NO + H2O}$.

\textbf{2.} É preciso ácido para remover, como água, o grupo hidróxi do intermediário de adição,
mas ácido demais protona a hidroxilamina, que então deixa de se adicionar.

\textbf{3.} Não: os dois carbonos ligados ao carbono C=N são equivalentes pela
simetria do anel.

\textbf{4.} Na oxima \textit{E}, o OH aponta para longe do carbono C2, que porta a metila (anti a
ele); na oxima \textit{Z}, na direção dele.

\textbf{5.} A inversão em um nitrogênio de oxima que porta um oxigênio eletronegativo tem uma barreira
alta.

\textbf{6.} Ele protona (ou sulfona) o grupo hidróxi, tornando-o um bom grupo de
saída, a água.

\textbf{7.} A ligação C--C do anel anti ao grupo de saída migra do carbono para o nitrogênio,
que é assim inserido no anel: seis átomos passam a sete.

\textbf{8.} Um íon nitrílio cíclico.

\textbf{9.} $\ce{C6H11NO -> C6H11NO}$: uma isomerização, de oxima a lactama.

\textbf{10.} Caprolactama, azepan-2-ona, a amida cíclica de sete membros do
ácido 6-amino-hexanoico; sua polimerização por abertura de anel dá o náilon-6.

\textbf{11.} \qty{98.0}{g/mol} e \qty{113.0}{g/mol}.

\textbf{12.} $\qty{1.0e6}{g}/\qty{98.0}{g/mol} = \qty{10204}{mol}$; $\times 0.98 \times 0.98 = \qty{9800}{mol}$; $\times \qty{113.0}{g/mol} = \qty{1.11e6}{g}$, \qty{1.11}{t}.

\textbf{13.} Oxima: $\num{10204} \times 0.98 = \qty{1.00e4}{mol}$; $\ce{H2SO4}$: $\qty{1.30e4}{mol}$, dando outros tantos mols de
$\ce{(NH4)2SO4}$: $\qty{1.30e4}{mol} \times \qty{132.1}{g/mol} = \qty{1.7}{t}$, mais que a caprolactama.

\textbf{14.} $\ce{C6H10O + RCO3H -> C6H10O2 + RCOOH}$.

\textbf{15.} O aduto tetraédrico do perácido sobre o carbono carbonílico (o
intermediário de Criegee), um peroxi-hemicetal.

\textbf{16.} Um grupo $\ce{CH2}$ do anel (os dois são equivalentes) migra para o oxigênio
peroxídico mais próximo; o carboxilato $\ce{RCOO-}$ sai.

\textbf{17.} $\varepsilon$-Caprolactona, oxepan-2-ona; sua polimerização por abertura de anel dá um
poliéster biodegradável, a poli($\varepsilon$-caprolactona).

\textbf{18.} Seu único subproduto é a água, enquanto um perácido deixa um mol de
ácido carboxílico por mol de lactona e é ele próprio perigoso de armazenar.

\textbf{19.} O carbono secundário C2 migra: 7-metiloxepan-2-ona.

\textbf{20.} Sim: o carbono que migra conserva sua configuração.

\textbf{21.} 7-Metilazepan-2-ona (C2 anti, migra).

\textbf{22.} 3-Metilazepan-2-ona (C6 migra).

\textbf{23.} Beckmann: a geometria da oxima (o grupo anti migra).
Baeyer--Villiger: a \omterm{def:b3:radicals-carbenes:migratory}{aptidão migratória} (o carbono mais substituído migra).

\textbf{24.} Seus dois carbonos $\alpha$ são equivalentes: uma oxima, uma lactama, uma lactona.

\textbf{25.} Uma tonelada de ciclo-hexanona dá cerca de \qty{1.11}{t} de caprolactama.
\end{solution}
